\documentclass[10pt,a4paper]{amsart}
\usepackage[margin=19mm]{geometry}
\usepackage{amsmath,amssymb,mathtools}
\usepackage[scr=rsfs,cal=euler]{mathalfa}
\usepackage{xfrac}
\usepackage[unicode,hidelinks,bookmarks=false]{hyperref}
\newcommand{\ie}{i.e.\ }
\newcommand{\cf}{cf.\ }
\newcommand{\resp}{respectively\ }
\newcommand{\Subsection}{\subsection}
\providecommand{\nobreakdash}{\nobreak\hbox{-}}
\setlength{\emergencystretch}{2em}
\multlinegap0pt
\usepackage{xcolor}
\usepackage[normalem]{ulem}
\usepackage{amssymb}
\usepackage{mathtools}
\usepackage{bbm}
\usepackage{algorithm}\usepackage{algpseudocode}
\usepackage{tikz}
\usepackage{caption}
\usepackage{tcolorbox}
\newtheorem{lemma}{Lemma}
\newcommand{\R}{\mathbb{R}}
\newcommand{\Xalpha}{{X^\alpha}}
\newcommand{\diff}{\mathrm{d}}
\newcommand{\dv}{\mathrm{div}}
\title{From Consensus-Based Optimization to Evolution Strategies: Proof of Global Convergence}
\author{Massimo Fornasier, Hui Huang, Jona Klemenc, Greta Malaspina}
\date{Source: arXiv:2602.11677v1 (CC BY 4.0)}
\begin{document}
\maketitle

\noindent\emph{Source and licence.} From Consensus-Based Optimization to Evolution Strategies: Proof of Global Convergence. Version arXiv:2602.11677v1 (CC BY 4.0), distributed by arXiv under the Creative Commons Attribution 4.0 International licence, http://creativecommons.org/licenses/by/4.0/, as recorded in the arXiv licence metadata for this exact version and shown on the abstract page https://arxiv.org/abs/2602.11677v1. Full licence notice: \texttt{licenses/arxiv-2602.11677-CC-BY-4.0.txt}.\par\smallskip

\noindent\textit{Editorial scope.} Counted unit: the article's lemma lemma:ball\_measure (source0.tex lines 573--588). The article's complete original proof of it is delivered with it (source0.tex lines 1956--2046, one \texttt{proof} environment of 91 lines, which the article places away from the statement). No mathematics is written, repaired or re-derived: the statement and the proof are the article's own lines, in the article's own notation, and no retained line is substituted or re-flowed.  Retained uncounted as the source-local context the proof needs: lines 466--468.  Nothing was omitted from any retained span, so the package carries no omission record.  No retained line contains a citation, so the wrapper carries no bibliography and no bibliography entry.  No retained line contains a figure environment, an image-inclusion command or drawing code, so no image is delivered and the package declares no asset dependency.  Editorial formatting only, in addition: an AMS \texttt{amsart} preamble that restates the article's own declarations and macros used by the retained lines, namely the environments \texttt{lemma} on separate counters, together with the standard packages those lines need.  The retained environments print in this document as Lemma 1 in order of appearance, whereas the article prints the same items with the numbering of its own full document.  The complete native source of the article as distributed in the arXiv e-print for this exact version is bundled unchanged as \texttt{sources/194484/source0.tex}.
\begin{equation}\label{eq:FP_CBO_delta}
\partial_t\rho^\alpha_t(x) = \lambda \dv\left(\rho^\alpha_t(x)(x-\Xalpha(\rho^\alpha_t))\right) +\frac{\delta^2}{2}\Delta \rho^\alpha_t(x).
\end{equation}

\begin{lemma}\label{lemma:ball_measure}
Let {$\rho^\alpha$ % \in \mathcal{C}([0,T]; \mathcal{P}_2(\mathbb{R}^d))$ 
be a %weak 
smooth solution of equation \eqref{eq:FP_CBO_delta}} with fixed $\alpha,\lambda,\delta>0$, for a given $T>0$ and a given initial distribution $\rho_0$. Given any $r>0$ let us denote with $\phi_r:\R^d\longrightarrow[0,1]$ the function 
$$\phi_r(x) = \begin{cases}
    \exp\left(1-\frac{r^2}{r^2- |x-x^*|^2}\right) & \text{ if }  |x-x^*|<r\\
    0 & \text{ otherwise .}
\end{cases}$$
Then, for every $t\geq 0$ 
    \begin{equation}
        \rho^\alpha_t(B_r(x^*))\geq e^{-pt}\int_{\mathbb{R}^d}\phi_r(x)\rho_0(x)\diff x
    \end{equation}

with $$p =\max\left\{2\frac{\lambda c^{1/2}(c^{1/2}r+B)}{r(1-c)^2}+\frac{\delta(1+c^2)d}{r^2(1-c)^4}, \frac{4\lambda^2r(cr^2+B)}{(2c-1)\delta^2}\right\}$$
for any $B,c>0$ such that $\sup_{t\in[0,T]}|X^\alpha(\rho^\alpha_t)-x^*|\leq B$, and $(2c-1)c\geq d(1-c)^2.$
\end{lemma}

\begin{proof}[Proof of Lemma \ref{lemma:ball_measure}]
    % The proof is basically the same as that of Proposition 4.6 in \url{https://arxiv.org/pdf/2103.15130}. The only difference is that at the beginning 
    % $$T_2 = \frac{\delta^2}{2}\Delta\phi_r(x)$$
    % $$K_2 = \{x\ | \ -\lambda(x-\Xalpha(\rho_t))^\top(x-x^*)(r^2- |x-x^*|^2)^2>\tilde c r^2 \frac{\delta^2}{2} |x-x^*|^2\}.$$
Since $\phi_r(x)\in[0,1]$ and $\phi_r(x)=0$ whenever $ |x-x^*|\geq r$ we have that 
$$\rho^\alpha_t(B_r(x^*)) \geq \int_{\mathbb{R}^d}\phi_r(x) \diff\rho^\alpha_t(x).$$
Let us denote with $T_1(x)$ and $T_2(x)$ the following two quantities
\begin{equation*}
    \begin{aligned}
        &T_1(x) = -\lambda(x-\Xalpha(\rho^\alpha_t))^\top\nabla\phi_r(x)\\
        &T_2(x) = \frac{\delta^2}{2}\Delta\phi_r(x).
    \end{aligned}
\end{equation*}
By \eqref{eq:FP_CBO_delta} we have 
\begin{equation}\label{eq:T1T2}
    \frac{\diff}{\diff t}\int\phi_r(x)\diff\rho^\alpha_t(x) = \int\partial_t\rho^\alpha_t(x)\phi_r(x)\diff x = \int_{\Omega_r} T_1(x)+T_2(x)\diff\rho^\alpha_t(x),
\end{equation}
where the gradient and the Laplacian of $\phi_r$ are given by 
    $$
\nabla \phi_r(v) = -2r^2 \frac{v - v^*}{\left( r^2 - |v - v^*|_2^2 \right)^2} \phi_r(v),
$$
$$
\Delta \phi_r(v) = 2r^2 \left( \frac{2 \left( 2|v - v^*|_2^2 - r^2 \right)|v - v^*|_2^2 - d \left( r^2 - |v - v^*|_2^2 \right)^2}{\left( r^2 - |v - v^*|_2^2 \right)^4} \right) \phi_r(v)
$$
respectively, and $\Omega_r= \{x\in\R^d\ |\  |x-x^*|<r\}$. Introducing the following two subsets 
\begin{equation*}
    \begin{aligned}
        &A_1 = \{x\in\R^d\ | \  |x-x^*|> rc^{1/2}\}\\
        &A_2 = \left\{x\in\R^d\ | \ -\lambda(x-\Xalpha(\rho^\alpha_t))^\top(x-x^*)\left(r^2- |x-x^*|^2\right)^2>\frac{\delta^2}{2}(2c-1)r^2 |x-x^*|^2\right\}
    \end{aligned}
\end{equation*}
the integral in \eqref{eq:T1T2} can be rewritten as
\begin{equation}\begin{aligned}
    \frac{\diff}{\diff t}\int\phi_r(x)\diff\rho^\alpha_t(x) &=
    \int_{\Omega_r\cap A_1^c} T_1(x)+T_2(x)\diff\rho^\alpha_t(x) +
    \int_{\Omega_r\cap A_1\cap A_2^c} T_1(x)+T_2(x)\diff\rho^\alpha_t(x) \\& +
    \int_{\Omega_r\cap A_1\cap A_2} T_1(x)+T_2(x)\diff\rho^\alpha_t(x).  
\end{aligned}\end{equation}
We now bound the three integrals.\\
Let us first consider the case $x\in A_1^c\cap\Omega_r$, which is equivalent to $ |x-x^*|\leq rc^{1/2}.$ In this case we have 
$|x-\Xalpha(\rho^\alpha_t)|\leq |x-x^*|+|\Xalpha(\rho^\alpha_t)-x^*|\leq c^{1/2}r+B$
and $r^2- |x-x^*|^2\geq r^2(1-c).$ Using these inequalities in the definition of $T_1$ we get
\begin{equation}
    \begin{aligned}
        T_1(x) &= \frac{2\lambda r^2\phi_r(x)(x-\Xalpha(\rho^\alpha_t))^\top(x-x^*)}{\left(r^2- |x-x^*|^2\right)^2}\geq-\frac{2\lambda r^2\phi_r(x)|x-\Xalpha(\rho^\alpha_t)| |x-x^*|}{\left(r^2- |x-x^*|^2\right)^2}\\
        &\geq-2\frac{\lambda c^{1/2}(c^{1/2}r+B)}{r(1-c)^2}\phi_r(x) = -p_1\phi_r(x).\\
    \end{aligned}
\end{equation}
On the other hand, by definition of $T_2$ we have
\begin{equation}
    \begin{aligned}
        T_2(x) &= -\frac{\delta^2r^2\phi_r(x)}{\left(r^2- |x-x^*|^2\right)^4}\left(d\left(r^2- |x-x^*|^2\right)^2+2 |x-x^*|^2{\left(r^2-2 |x-x^*|^2\right)^2}\right)\\
        &\geq -\frac{\delta(1+c^2)d}{r^2(1-c)^4}\phi_r(x) = -p_2\phi_r(x).\\
    \end{aligned}
\end{equation}
We now prove that if $x\in \Omega_r\cap A_1\cap A_2^c$ then $T_1(x)+T_2(x)\geq 0.$ This is true if and only if 
\begin{equation}\label{eq:A2_T1T2}
    \left(-\lambda(x-\Xalpha(\rho^\alpha_t))^\top(x-x^*)+\frac{d\delta^2}{2}\right)\left(r^2- |x-x^*|^2\right)^2 \leq \delta^2 |x-x^*|^2\left(2 |x-x^*|^2-r^2\right).
\end{equation}
Since $x\in A_2^c$, $x\in A_1$ and the fact that $(1-c)^2d\leq(2c-1)c$ we have 
\begin{equation*}\begin{aligned}
    \left(-\lambda(x-\Xalpha(\rho^\alpha_t))^\top(x-x^*)+\frac{d\delta^2}{2}\right)\left(r^2- |x-x^*|^2\right)^2 &\leq \frac{\delta^2}{2} |x-x^*|^2(2c-1)r^2 + \frac{\delta^2}{2}(2c-1)cr^4\\
    &\leq \delta^2 |x-x^*|^2\left(2 |x-x^*|^2-r^2\right),
\end{aligned}\end{equation*}
which yields \eqref{eq:A2_T1T2}, and thus $T_1(x)+T_2(x)\geq 0$ for $x\in \Omega_r\cap A_1\cap A_2^c$.
Let us now consider the case $x\in \Omega_r\cap A_1\cap A_2$. By definition of $A_2$ and Cauchy Schwartz inequality we have 
\begin{equation}\begin{aligned}
    T_1(x)&\geq-\frac{2\lambda r^2\phi_r(x)}{(r^2- |x-x^*|^2)^2}|x-\Xalpha(\rho^\alpha_t)| |x-x^*|\geq \frac{4\lambda^2|x-\Xalpha(\rho^\alpha_t)|}{(2c-1)\delta^2}(x-\Xalpha(\rho^\alpha_t))^\top(x-x^*)\\
    & \geq -\frac{4\lambda^2r(cr^2+B)}{(2c-1)\delta^2}\phi_r(x) = -p_3\phi_r(x).
\end{aligned}\end{equation}
By definition of $T_2$ and $\Delta\phi_r(x)$ we have that $T_2(x)\geq 0 $ if 
\begin{equation}\label{eq:A3T2}
    2 |x-x^*|^2(2 |x-x^*|^2-r^2)\geq d(r^2-|x-x|^2)^2.
\end{equation}
Since $x\in\Omega_r\mathcal{A}_1\cap \mathcal{A}_2$ we have 
\begin{equation*}
  2 |x-x^*|^2(2 |x-x^*|^2-r^2)\geq 2r^4c(2c-1)\geq dr^4(1-c)^2\geq d(r^2-|x-x|^2)^2.
\end{equation*}
Thus \eqref{eq:A3T2} holds and $T_2(x)\geq 0$.
Putting everything together, we then get 
\begin{equation}\begin{aligned}
    \frac{\diff}{\diff t}\int\phi_r(x)\diff\rho^\alpha_t(x) &=
    \int_{\Omega_r\cap A_1^c} T_1(x)+T_2(x)\diff\rho^\alpha_t(x) +
    \int_{\Omega_r\cap A_1\cap A_2^c} T_1(x)+T_2(x)\diff\rho^\alpha_t(x) \\& +
    \int_{\Omega_r\cap A_1\cap A_2} T_1(x)+T_2(x)\diff\rho^\alpha_t(x)\\
    &\geq -(p_1+p_2)\int_{\Omega_r\cap A_1^c} \phi_r(x)\diff\rho^\alpha_t(x) 
    -p_3\int_{\Omega_r\cap A_1\cap A_2} \phi_r(x)\diff\rho^\alpha_t(x)\\
    &\geq -p\int\phi_r(x)\diff\rho^\alpha_t(x),
\end{aligned}\end{equation}
 with $p=\max\{p_1+p_2,p_3\}$. The thesis then follows from Gronwall's inequality.
\end{proof}

\end{document}
